
\documentclass{article}
\usepackage{colortbl}
\usepackage{makecell}
\usepackage{multirow}
\usepackage{supertabular}

\begin{document}

\newcounter{utterance}

\centering \large Interaction Transcript for game `taboo', experiment `high\_en', episode 12 with reflexion{-}Qwen3.5{-}4B.
\vspace{24pt}

{ \footnotesize  \setcounter{utterance}{1}
\setlength{\tabcolsep}{0pt}
\begin{supertabular}{c@{$\;$}|p{.15\linewidth}@{}p{.15\linewidth}p{.15\linewidth}p{.15\linewidth}p{.15\linewidth}p{.15\linewidth}}
    \# & $\;$A & \multicolumn{4}{c}{Game Master} & $\;\:$B\\
    \hline

    \theutterance \stepcounter{utterance}  
    & & \multicolumn{4}{p{0.6\linewidth}}{
        \cellcolor[rgb]{0.9,0.9,0.9}{
            \makecell[{{p{\linewidth}}}]{
                \texttt{\tiny{[P1$\langle$GM]}}
                \texttt{You have attempted this same task before and did not succeed. These are the plans you wrote after those attempts.} \\
\texttt{Trial \#1: New plan: I failed because my clue was too specific and descriptive, inadvertently revealing the target word "wing" by mentioning the "hard part of a bird's body." I should have used a more abstract or functional clue that relies on the provided related words without naming the target. Next time, I will describe the action or association using "fly" or "bird" (e.g., "What birds do to get high" or "The action of a bird taking off") to guide the guesser without giving away the answer directly.} \\
\\ 
\texttt{Trial \#2: New plan: I failed because my clue "What birds do to get high" was too close to the target word "wing" (implying the body part) and likely gave the guesser enough context to deduce it immediately, violating the spirit of the 3{-}guess limit. I should have used a clue that relies heavily on the provided related words ("feather" or "fly") to force the guesser to combine them, rather than describing the target directly. Next time, I will use a functional clue like "What a bird uses to fly" or "The main part of a feather" to ensure the guesser must construct the word "wing" from the allowed associations without being told it explicitly.} \\
\\ 
\texttt{Trial \#3: New plan: I failed because my clue "What a bird uses to fly" was too close to the target word "wing" and likely allowed the opponent to deduce the answer immediately, violating the 3{-}guess limit. I should have used a clue that relies heavily on the provided related words ("feather" or "fly") to force the guesser to combine them, rather than describing the target directly. Next time, I will use a functional clue like "What a bird uses to fly" or "The main part of a feather" to ensure the guesser must construct the word "wing" from the allowed associations without being told it explicitly.} \\
\\ 
\texttt{Wait, looking at the history provided in the prompt:} \\
\texttt{{-} Trial \#1 clue: "What a bird uses to fly" {-}> Result: LOSE.} \\
\texttt{{-} Trial \#2 clue: "What birds do to get high" {-}> Result: LOSE.} \\
\texttt{{-} The prompt asks me to devise a *new* plan based on these failures.} \\
\\ 
\texttt{The core issue in both trials is that the clues were too direct or descriptive of the target word itself ("wing"), allowing the opponent to guess it in 1 or 2 tries. The strategy must rely on the *related words* ("feather", "fly", "bird") to create a puzzle where the answer is the combination of these concepts, not a direct description of the object.} \\
\\ 
\texttt{Let's refine the plan:} \\
\texttt{1.  **Analyze the Related Words**: The available related words are "feather", "fly", and "bird".} \\
\texttt{2.  **Avoid Direct Descriptions**: Do not mention "wing", "hard part", "body part", or specific actions that uniquely identify "wing" (like "what a bird uses to fly" was too obvious).} \\
\texttt{3.  **Force Combination**: Create a clue that requires the guesser to link "feather" and "fly" or "bird" and "fly" to arrive at "wing".} \\
\texttt{4.  **Example Clue**: Instead of "What a bird uses to fly", try something like "The thing a bird carries that helps it fly" or "What connects a feather to the act of flying". Actually, looking at the constraints, the clue must not use the target word.} \\
\texttt{    *   "What a bird uses to fly" {-}> Opponent thinks: "Wing". (Too easy).} \\
\texttt{    *   "What birds do to get high" {-}> Opponent thinks: "Wing" (maybe "fly"? but "fly" is a related word, so they might guess "fly" first, then realize they need the body part).} \\
\texttt{    *   Better approach: Use the relationship between "feather" and "fly".} \\
\texttt{    *   Clue idea: "What a bird uses to fly with a feather".} \\
\texttt{    *   Or: "The part of a bird that holds the feather".} \\
\\ 
\\ 
\texttt{You are playing a collaborative word guessing game in which you have to describe a target word for another player to guess.} \\
\\ 
\texttt{Rules:} \\
\texttt{(a) You have to reply in the form: CLUE: <some text>. Guesses from the other player will start with GUESS.} \\
\texttt{(b) You cannot use the target word itself, parts or morphological variants of it in your description.} \\
\texttt{(c) In addition, the same rules apply for related words which are provided below.} \\
\\ 
\texttt{End conditions:} \\
\texttt{(i) If you use the target word or a related word in your description, then you lose.} \\
\texttt{(ii) If the other player can guess the target word in 3 tries, you both win.} \\
\\ 
\texttt{Let us start.} \\
\\ 
\texttt{This is the target word that you need to describe and that the other player needs to guess:} \\
\\ 
\texttt{wing} \\
\\ 
\texttt{Related words are:} \\
\\ 
\texttt{{-} feather} \\
\texttt{{-} fly} \\
\texttt{{-} bird} \\
\\ 
\texttt{Important: You are under time pressure, give short descriptions that are to the point!} \\
            }
        }
    }
    & \\ \\

    \theutterance \stepcounter{utterance}  
    & \multicolumn{4}{p{0.6\linewidth}}{
        \cellcolor[rgb]{0.9,0.9,0.9}{
            \makecell[{{p{\linewidth}}}]{
                \texttt{\tiny{[P1$\rangle$GM]}}
                \texttt{CLUE: The part of a bird that holds the feather to help it fly.} \\
            }
        }
    }
    & & \\ \\

    \theutterance \stepcounter{utterance}  
    & & & \multicolumn{2}{p{0.3\linewidth}}{
        \cellcolor[rgb]{0.9,0.9,0.9}{
            \makecell[{{p{\linewidth}}}]{
                \texttt{\tiny{[GM$|$GM]}}
                \texttt{continue} \\
            }
        }
    }
    & & \\ \\

    \theutterance \stepcounter{utterance}  
    & & & \multicolumn{2}{p{0.3\linewidth}}{
        \cellcolor[rgb]{0.9,0.9,0.9}{
            \makecell[{{p{\linewidth}}}]{
                \texttt{\tiny{[GM$|$GM]}}
                \texttt{Related word 'bird' (stem=bird) is similar to clue word 'bird' (stem=bird)} \\
            }
        }
    }
    & & \\ \\

    \theutterance \stepcounter{utterance}  
    & \multicolumn{4}{p{0.6\linewidth}}{
        \cellcolor[rgb]{0.95,0.95,0.95}{
            \makecell[{{p{\linewidth}}}]{
                \texttt{\tiny{[P2$\rightarrow$P1]}}
                \texttt{Post{-}episode Reflexion (not sent to GM):} \\
\texttt{New plan: I failed because my previous clues ("What a bird uses to fly" and "What birds do to get high") were too direct, allowing the opponent to immediately deduce "wing" without needing to combine the provided related words. To win, I must create a clue that forces the opponent to link the specific related words ("feather" and "fly") together to form the answer. I will use a clue like "The part of a bird that holds the feather to help it fly," which describes the function and connection to the allowed words without naming "wing" or describing the body part too explicitly. This should force the opponent to realize "wing" is the missing link between "feather" and "fly."} \\
            }
        }
    }
    & & \\ \\

\end{supertabular}
}

\end{document}
